\subsection{Проектирование модулей}

Стоит учитывать, что целью работы является создание программного модуля распределения вершин, а не графовой БД, то есть программа, создаваемая в рамках работы, не должна обладать всеми функциями, ожидаемые от базы данных, а только имитировать те, что нужны для работы и тестирования составляющих модуля.

Для выполнения поставленных задач программа должна состоять из следующих модулей:
\begin{itemize}
    \item мастер;
    \item хранилище;
    \item имитация шины.
\end{itemize}

Мастер при этом состоит из:
\begin{itemize}
    \item статического оптимизатора;
    \item модуля управления запросами.
\end{itemize}

\textit{Мастер}

Мастер-модуль предназначен для централизованного управления системой,
координации взаимодействия между остальными модулями и выполнения операций
инициализации и обслуживания распределённого графового хранилища, а также служить входной точкой в систему.

Основные функции модуля:
\begin{itemize}
    \item загрузка конфигурации системы и создание объектов модулей;
    \item управление жизненным циклом хранилищ и компонентов системы;
    \item загрузка графовых данных в систему;
    \item запуск процедур оптимизации распределения данных;
    \item инициализация выполнения запросов;
    \item вывод результата запросов.
\end{itemize}

Следовательно методами модуля должны быть:
\begin{itemize}
    \item подключить шину;
    \item отключить шину;
    \item подключить хранилище;
    \item отключить хранилище;
    \item добавить вершину в хранилище (с исходящими рёбрами);
    \item добавить вершину (с исходящими рёбрами);
    \item запуск статической оптимизации;
    \item отправка запроса на поиск пути.
\end{itemize}

\textit{Модуль статического оптимизатора}

Модуль статического оптимизатора является подмодулем и фактической составляющей мастера и предназначен для построения и улучшения распределения графовых данных между шардами на основе структуры графа и накопленной статистики запросов.

Основные функции модуля:

\begin{itemize}
    \item балансировка размеров партиций;
    \item применение алгоритма Кернигана-Лина для улучшения распределения.
\end{itemize}

Следовательно у модуля будет один метод: провести оптимизацию распределения пары хранилищ.

\textit{Модуль управления запросами}

Модуль управления запросами предназначен является подмодулем и фактической составляющей мастера и предназначен для управления выполнением запросами.

Основные функции модуля:

\begin{itemize}
    \item управление выполнением запросов на поиск пути;
    \item сбор статистики выполнения сформированных запросов;
    \item подготовка данных для анализа эффективности оптимизации.
\end{itemize}

Следовательно у модуля будет один метод: выполнить запрос на поиск пути.

\textit{Модуль хранилища}

Модуль хранилища реализует хранение графовых данных, выполняет запросы на поиск пути.

Основные функции модуля:

\begin{itemize}
    \item хранение вершин и рёбер графа;
    \item добавление вершин и рёбер графа;
    \item расчёт объективной функции добавления вершины;
    \item выполнение локального поиска пути внутри хранилища;
    \item обработка запросов на чтение данных;
    \item операции перемещения вершины между хранилищами.
\end{itemize}


Следовательно методами модуля должны быть:
\begin{itemize}
    \item подключиться к шине;
    \item отключиться от шины;
    \item добавить вершину (с исходящими рёбрами);
    \item расчитать объективную функцию добавления вершины;
    \item отправить вершину в другое хранилище;
    \item выполнить запрос на поиск пути.
\end{itemize}

\textit{Модуль имитации шины}

Модуль имитации шины предназначен для моделирования сетевого взаимодействия между распределёнными узлами хранения и мастером.

Основной функцией модуля является быть прослойкой для всех взаимодействий между мастером и хранилищами.

Следовательно методами модуля должны быть:
\begin{itemize}
    \item подключить хранилище;
    \item отключить  хранилище;
    \item разослать сообщения.
    
\end{itemize}

В рассылку сообщений входят запросы и ответы:
\begin{itemize}
    \item о вычислении объективной функции новой вершины;
    \item о добавлении вершины;
    \item о запросе данных вершины;
    \item о перемещении вершины;
    \item о запросе граничных вершин.
\end{itemize}
